% =============================================================================
\chapter{Introdução}\label{cap:introducao}
% =============================================================================

O Projeto Analítico A3 é a avaliação prática e continuada da unidade curricular
Análise de Dados e Big Data. Segundo as diretrizes da disciplina
\cite{usjt2026}, cada grupo deve diagnosticar um problema real, coletar e
sanitizar uma base de dados pública por meio de um processo de extração,
transformação e carga (ETL), realizar análises exploratórias e preditivas e
comunicar os resultados por meio de \textit{Data Storytelling} e visualização,
respeitando a ética e a Lei Geral de Proteção de Dados Pessoais (LGPD).

Este documento consolida, em um único volume, tudo o que o grupo produziu nos
três \textit{checkpoints} da A3: o diagnóstico e o escopo, a auditoria de
qualidade e o ETL, a análise exploratória, a modelagem preditiva, o painel
interativo, a narrativa executiva, a reflexão ética e as instruções de
reprodução. Ele foi escrito a partir do estado do repositório do projeto na
revisão \code{9aaf3fb} (novembro de 2026) e das saídas registradas nos
notebooks executados.

\section{Contexto e cenário}\label{sec:contexto}

O grupo assume o papel do time de dados de uma \textbf{mesa de execução de
criptoativos}, área que recebe ordens grandes de clientes institucionais
(fundos, tesourarias de empresas e \textit{family offices}) e precisa
executá-las no mercado à vista de BTC/USDT da Binance, a corretora de maior
liquidez do mundo. Executar mal custa caro: uma ordem grande colocada no
momento errado consome vários níveis de preço do livro de ofertas, paga
\textit{slippage} (diferença entre o preço esperado e o efetivamente obtido) e
fica exposta a movimentos bruscos.

O período escolhido, de 19 a 25 de fevereiro de 2025, foi uma semana de forte
estresse. Em 21 de fevereiro a corretora Bybit sofreu o maior roubo de
criptoativos já registrado, de aproximadamente \usd~1,5 bilhão em ether (ETH),
atribuído pelo FBI a agentes norte-coreanos \cite{fbi2025}. Nos dias seguintes,
o bitcoin, que havia tocado \usd~99.475 em 21 de fevereiro, caiu até
\usd~86.051 em 25 de fevereiro, uma queda de 13,5\,\%. Somente em 25 de
fevereiro ocorreram 10,3 milhões de negócios, 37\,\% de toda a semana.

\section{Problema e objetivos}\label{sec:objetivos}

A mesa precisa de evidências, e não de intuição, para decidir \textbf{quando}
e \textbf{como} executar ordens e para \textbf{calibrar seus controles de risco}
em dias de estresse. Dados de \textit{tick}, isto é, cada negócio individual com
precisão de milissegundo, permitem medir liquidez, volatilidade e pressão
compradora ou vendedora em uma granularidade que os candles diários escondem.

O \textbf{objetivo geral} do projeto é transformar uma base pública de
27,9 milhões de negócios em evidências confiáveis para as decisões de uma mesa
de execução. Os \textbf{objetivos específicos} são:

\begin{alineas}
  \item auditar a qualidade da base bruta e documentar cada problema encontrado,
        aplicando o princípio \textit{Garbage In, Garbage Out} (GIGO);
  \item construir um pipeline de ETL reprodutível, com verificação de
        integridade, contrato de qualidade e validação cruzada contra a fonte
        primária;
  \item responder, com métricas e agrupamentos, a três perguntas de negócio
        sobre liquidez, estresse e fluxo de ordens;
  \item avaliar, sem vazamento de informação futura, se é possível prever a
        volatilidade e a direção do preço no intervalo seguinte;
  \item comunicar os achados em um painel interativo e em uma narrativa
        orientada à decisão;
  \item analisar criticamente a ética, a conformidade com a LGPD, a licença da
        fonte e as limitações do estudo.
\end{alineas}

\section{Perguntas de negócio}\label{sec:perguntas}

O \autoref{qua:perguntas} apresenta as três perguntas de negócio definidas no
\textit{Checkpoint}~1, com a decisão que cada uma apoia, as métricas usadas e a
hipótese inicial do grupo.

\begin{quadro}[htbp]
\caption{\label{qua:perguntas}Perguntas de negócio, decisões apoiadas e hipóteses}
\begin{corpotabela}
\begin{tabularx}{\textwidth}{|C{0.6cm}|Y|L{3.3cm}|L{3.6cm}|L{3.2cm}|}
\hline
\textbf{Id} & \textbf{Pergunta} & \textbf{Decisão apoiada} & \textbf{Métricas} & \textbf{Hipótese} \\ \hline
P1 & Em quais horários do dia (UTC) o BTC/USDT concentra maior volume e número
de negócios, e como a volatilidade acompanha essa liquidez? &
Janela em que a mesa deve executar ou fatiar ordens grandes para reduzir o
impacto de preço. &
Volume em dólares e número de negócios por minuto (média e mediana por hora);
retorno absoluto e amplitude por hora; mapa de calor dia $\times$ hora. &
A liquidez se concentra na abertura do mercado americano, mas a volatilidade
também, o que cria um \textit{trade-off}. \\ \hline
P2 & Como volume, volatilidade e perfil de tamanho dos negócios mudaram nos
dias de estresse em relação aos dias calmos? Quem concentra o volume? &
Calibração de limites de risco e alertas; dimensionamento da capacidade
operacional. &
Amplitude diária, volatilidade realizada, volume, número de negócios;
participação de cada faixa de tamanho (\code{groupby}). &
Em estresse o volume se multiplica e poucos negócios grandes concentram a
maior parte do valor. \\ \hline
P3 & O desequilíbrio entre compras e vendas agressoras (OFI) está associado ao
retorno do mesmo minuto? Ele antecipa o retorno do minuto seguinte? &
Viabilidade de um sinal ou modelo preditivo de curto prazo. &
OFI por minuto; correlações de Pearson e Spearman com defasagens de 0 a 10
minutos; saldo comprador diário. &
Forte associação no mesmo minuto, mas pouco ou nenhum poder preditivo
(eficiência de mercado). \\ \hline
\end{tabularx}
\end{corpotabela}
\fonteautores[Elaborado]
\end{quadro}

\section{Equipe e cronograma de entregas}\label{sec:equipe}

O grupo é formado por dois integrantes, dentro do limite de dois a três alunos
estabelecido pelas diretrizes da A3: Gabriel Almeida Portela (RA 825233281) e
Daniel Almeida Portela (RA 825234443). A A3 vale 40 pontos e é entregue em três
\textit{checkpoints}, resumidos no \autoref{qua:checkpoints} com a situação de
cada entrega no momento desta documentação.

\begin{quadro}[htbp]
\caption{\label{qua:checkpoints}Cronograma de entregas da A3 e situação do projeto}
\begin{corpotabela}
\begin{tabularx}{\textwidth}{|L{3.0cm}|C{1.5cm}|C{1.3cm}|Y|L{2.6cm}|}
\hline
\textbf{\textit{Checkpoint}} & \textbf{Prazo} & \textbf{Valor} & \textbf{Entregáveis do projeto} & \textbf{Situação} \\ \hline
1 -- Diagnóstico e escopo & Semana 5 & 5,0 pts &
\code{CHECKPOINT_1.md}: grupo, cenário, base, dicionário de dados e três
perguntas de negócio. & Concluído \\ \hline
2 -- Auditoria e ETL & Semana 10 & 15,0 pts &
\code{notebooks/A3_CP2_auditoria_etl_eda.ipynb}: auditoria, limpeza, validação
cruzada, carga e EDA inicial. & Concluído \\ \hline
3 -- Entrega final e defesa & Semana 18 & 20,0 pts &
Código consolidado (\code{src/}), notebook de modelagem e
\textit{storytelling}, painel interativo (\code{app/}), análise de ética e LGPD
e esta documentação. & Concluído, exceto os \textit{slides} da defesa \\ \hline
\end{tabularx}
\end{corpotabela}
\fonte{Elaborado pelos autores (2026), com base em \citeonline{usjt2026}.}
\end{quadro}

\section{Organização do documento}\label{sec:organizacao}

O \autoref{cap:fundamentacao} apresenta os conceitos de microestrutura de
mercado, qualidade de dados e validação temporal que sustentam o trabalho. O
\autoref{cap:base} descreve a base de dados e o dicionário de dados. O
\autoref{cap:arquitetura} mostra a arquitetura da solução e a organização do
código. O \autoref{cap:etl} detalha a auditoria de qualidade e o ETL. O
\autoref{cap:eda} reúne a análise exploratória e as respostas às perguntas de
negócio. O \autoref{cap:modelagem} trata da modelagem preditiva. O
\autoref{cap:dashboard} documenta o painel interativo e o
\autoref{cap:storytelling}, a narrativa executiva. O \autoref{cap:etica}
discute ética, LGPD e limitações; o \autoref{cap:reprodutibilidade}, testes e
reprodutibilidade; e o \autoref{cap:consideracoes} apresenta as considerações
finais. Os apêndices trazem a matriz de atendimento aos requisitos da A3, o
dicionário dos produtos analíticos e a lista completa de atributos dos
modelos.
